\documentclass[../main.tex]{subfiles}
\begin{document}

% Nguồn số liệu của chương: repo SUMO (PLAN.md, config/*.yaml, rules/bounds/param_bounds.yaml,
% src/thesis_sumo/rules/schema.py, docs/PROGRESS.md đến R45, config/llm.yaml, config/site.proxy.yaml, report/2026-10-28, report/2026-10-29). Cập nhật 2026-10-31 (bản tạm thời);
% cập nhật 2026-11-09: report/2026-11-08/README.md, data/proxy/targets.yaml, prompts v003--v005.
% Các mục [TODO] cần bổ sung trước khi nộp.

Chương~\ref{chapter:Related_works} đã chỉ ra rằng chưa có công trình nào dùng LLM sinh cấu trúc luật hành vi tài xế cho dòng xe hỗn hợp, có kiểm tra hợp lệ và so sánh công bằng với một bộ tối ưu, và đã đặt ra yêu cầu cho hệ thống. Chương này trình bày phương pháp đáp ứng các yêu cầu đó. Một mô hình ngôn ngữ lớn (LLM) chạy cục bộ sinh và tinh chỉnh lặp bộ luật hành vi tài xế; trình mô phỏng SUMO~\cite{lopez2018sumo} và một bộ điều khiển TraCI thực thi các luật đó. Mục~\ref{section:3.1} nêu kiến trúc và trạng thái triển khai. Mục~\ref{section:3.2} mô tả kịch bản và mạng mô phỏng. Các mục~\ref{section:3.3} đến~\ref{section:3.5} mô tả cách biểu diễn luật, kiểm tra hợp lệ, biên dịch và thực thi. Mục~\ref{section:3.6} và~\ref{section:3.7} mô tả dữ liệu mục tiêu và chỉ số. Mục~\ref{section:3.8} mô tả vòng hiệu chỉnh và các phương án đối sánh.

\section{Tổng quan phương pháp}
\label{section:3.1}

Phương pháp chia hai việc cho hai bên (Hình~\ref{fig:1.1}). Việc \emph{sinh luật} thuộc về LLM: LLM nhận mô tả nút giao, từ điển tham số được phép chỉnh và phản hồi của vòng trước, rồi trả về một đặc tả luật (rule spec) dạng JSON. LLM chạy ngoài vòng mô phỏng, không điều khiển xe theo thời gian thực, và không sinh hay sửa mạng đường. Việc \emph{thực thi luật} thuộc về mã xác định: đặc tả được biên dịch thành (i) tham số tĩnh của kiểu phương tiện (vType), mà SUMO nạp một lần khi khởi động, và (ii) các luật có điều kiện, mà bộ điều khiển TraCI đánh giá ở mỗi bước mô phỏng.

Bảng~\ref{table:3.1} cho trạng thái triển khai của từng thành phần tại thời điểm viết.

\begin{table}[ht]
\centering
\caption{Trạng thái triển khai các thành phần của phương pháp.}
\label{table:3.1}
\begin{tabular}{|p{6.5cm}|p{7.5cm}|}
\hline
\textbf{Thành phần} & \textbf{Trạng thái} \\ \hline
Mạng tái lập được (đèn ba pha, làn rẽ phải liên tục ở 30~m cuối, kết nối rẽ trái trái phép, hai làn vỉa hè), nhu cầu, cấu hình, trình chạy mô phỏng & Đã triển khai, có kiểm thử tự động \\ \hline
Thí nghiệm khả thi Phase 3A--3D (chặn làn rẽ phải, vỉa hè, đi ngược chiều, xe chờ giữa nút) & Đã thực hiện; đi ngược chiều được hoãn \\ \hline
Lược đồ đặc tả luật, miền tham số, bộ kiểm tra hợp lệ, bộ biên dịch & Đã triển khai, có kiểm thử tự động \\ \hline
Bộ điều khiển TraCI với sáu luật có điều kiện & Đã triển khai, có kiểm thử tự động. Từ 04/11 luật \texttt{uturn\_bypass} thuộc cấu trúc B1 \\ \hline
Chỉ số M1--M10, cổng GEH (G0), hàm khoảng cách, điểm tổng & Đã triển khai, có kiểm thử tự động \\ \hline
Khung đánh giá chung và bộ đếm ngân sách & Đã triển khai, có kiểm thử tự động \\ \hline
Trình khách LLM, prompt, phản hồi, vòng lặp LLM (P1, P2) & Đã triển khai, có kiểm thử tự động. P1, P2 với \texttt{qwen3.5:9b} đã chạy trên mạng 04/11, không đặc tả nào đủ điều kiện (Mục~\ref{subsection:campaigns_n1108}) \\ \hline
Phương án Optuna (B2, H) & Đã triển khai. B2 xong 36/36 lần chạy trên mạng F2 và 24/24 lần chạy trên mỗi mạng ngày 28/10 (Mục~\ref{subsection:net_versions}), 36/36 lần chạy trên mạng 04/11; H chưa chạy \\ \hline
Đánh giá cuối trên ngày đánh giá (E1, E2) & Dự kiến \\ \hline
\end{tabular}
\end{table}

\section{Kịch bản và mạng mô phỏng}
\label{section:3.2}

\subsection{Phạm vi kịch bản}

Kịch bản là nút giao có đèn, bốn nhánh, Lê Văn Lương -- Hoàng Đạo Thúy (quận Thanh Xuân, Hà Nội), cao điểm sáng 07:30--09:00. Có hai loại phương tiện: ô tô (vClass \texttt{passenger}) và xe máy (vClass \texttt{motorcycle}). Xe vi phạm được tách thành lớp riêng: xe máy vi phạm dùng vClass \texttt{custom1}, ô tô vi phạm dùng \texttt{custom2}. Chỉ hai lớp này được vào các không gian mà xe tuân thủ không được vào. Người đi bộ, xe buýt, xe tải, đỗ xe dưới lòng đường và chọn lộ trình không nằm trong kịch bản.

Kịch bản gồm 11 hành vi ở Bảng~\ref{table:1.1}, chia thành ba nhóm khác nhau về cơ chế: F1 (không tuân thủ tín hiệu) gồm hành vi 1--3; F2 (tìm khoảng trống khi ùn) gồm hành vi 4--8, chung một cơ chế là hàng chờ dài thì tài xế tìm không gian khác để lên đầu hàng; và F3 (tương tác trong nút) gồm hành vi 9--11. Hành vi 8 (đi ngược chiều) đã qua thí nghiệm khả thi nhưng được hoãn, chưa có trong mô hình (Mục~\ref{subsection:phase3}). Mọi nhánh có làn rẽ phải liên tục, nên rẽ phải khi đèn đỏ không phải vi phạm và các chỉ số không tính nó là vi phạm. Hành vi 6 là việc xe có hướng đi thẳng, rẽ trái hoặc quay đầu đứng chờ đèn đỏ của chính nó trên lane 0 trong 30~m cuối trước vạch dừng. Đi thẳng hoặc rẽ trái trên làn này khi đèn xanh là đúng luật.

\subsection{Xây dựng mạng lưới}

Mạng được sinh hoàn toàn bằng script từ dữ liệu OpenStreetMap của khu vực nút (Hình~\ref{fig:3.net}). \texttt{netconvert} giữ các tuyến trong bán kính 300~m quanh tâm nút, chỉ giữ cạnh cho phép ô tô và xe máy, đặt kiểu đèn tĩnh, rồi bỏ các nút và đèn phụ trên các nhánh. Các chỉnh sửa sau đó (bề rộng làn, kết nối rẽ, chương trình đèn) đi qua các tệp vá plain XML sinh từ cấu hình hiện trường; không ai sửa tay \texttt{.net.xml}. Một lệnh duy nhất dựng lại toàn bộ mạng.

\begin{figure}[ht]
\centering
\includegraphics[width=0.9\textwidth]{Figure/net_diagram.png}
\caption{Mạng mô phỏng sinh từ OpenStreetMap với cấu hình đại diện. Bốn nhánh vào: LVL\_W và LVL\_E (Lê Văn Lương phía tây và phía đông), HDT\_N và HDT\_S (Hoàng Đạo Thúy phía bắc và phía nam). Đường đỏ: liên kết chỉ mở cho xe vi phạm (\texttt{custom1}, \texttt{custom2}); \texttt{711526326\#1.sidewalk} là làn vỉa hè. Thước tỷ lệ 100~m.}
\label{fig:3.net}
\end{figure}

Các hành vi cần không gian riêng dùng bốn bản vá tĩnh. Những kết nối và làn nằm ngoài phần đường dành cho xe tuân thủ chỉ mở cho lớp xe vi phạm. LLM không sinh hay sửa các bản vá này. Hình~\ref{fig:3.1} minh họa vị trí của chúng, Bảng~\ref{table:3.0} liệt kê chi tiết.

\begin{figure}[ht]
\centering
\begin{tikzpicture}[x=1cm, y=1cm, font=\footnotesize, >={Stealth[length=2.2mm]}]
  % đường
  \fill[gray!15] (-5.2,-0.95) rectangle (5.2,0.95);
  \fill[gray!15] (-0.95,-4.2) rectangle (0.95,4.2);
  \fill[gray!35] (-0.95,-0.95) rectangle (0.95,0.95);
  \draw[gray] (-5.2,-0.95) -- (-0.95,-0.95) -- (-0.95,-4.2);
  \draw[gray] (-5.2,0.95) -- (-0.95,0.95) -- (-0.95,4.2);
  \draw[gray] (5.2,-0.95) -- (0.95,-0.95) -- (0.95,-4.2);
  \draw[gray] (5.2,0.95) -- (0.95,0.95) -- (0.95,4.2);
  % làn rẽ phải liên tục: 30 m cuối của cả bốn nhánh
  \fill[green!35!black!25] (-1.9,-0.95) rectangle (-0.95,-0.6);
  \fill[green!35!black!25] (0.95,0.6) rectangle (1.9,0.95);
  \fill[green!35!black!25] (-0.95,0.95) rectangle (-0.6,1.9);
  \fill[green!35!black!25] (0.6,-1.9) rectangle (0.95,-0.95);
  % ranh giới đoạn 30 m cuối, nhánh tây
  \draw[very thick, dotted] (-1.9,-0.95) -- (-1.9,0);
  \draw[<->] (-1.9,0.12) -- node[above]{\tamthoi{30 m}} (-0.95,0.12);
  % hai làn vỉa hè
  \fill[orange!45] (-5.2,-1.35) rectangle (-0.95,-1.0);
  \fill[orange!45] (0.95,1.0) rectangle (5.2,1.35);
  % vạch dừng
  \draw[thick] (-0.95,-0.95) -- (-0.95,0);
  % xe vi phạm chờ giữa nút, từ nhánh Hoàng Đạo Thúy phía nam
  \draw[red, very thick, -{Stealth}] (0.3,-3.4) -- (0.3,-0.75) to[out=90,in=0] (-0.6,0.62) -- (-3.4,0.62);
  \foreach \p in {(0.27,-0.25),(0.12,0.2)} { \fill[red] \p circle (0.09); }
  % nhãn
  \node[anchor=east] at (-5.3,0) {Tây (Lê Văn Lương)};
  \node[anchor=west] at (5.3,0) {Đông};
  \node[anchor=south] at (0,4.3) {Bắc};
  \node[anchor=north] at (0,-4.3) {Nam};
  \node[green!35!black, anchor=north west, align=left] at (-5.2,-1.5) {làn rẽ phải liên tục (lane 0),\\30 m cuối của mỗi nhánh};
  \node[orange!60!black, anchor=north west, align=left] at (-5.2,-2.45) {làn vỉa hè, chỉ \texttt{custom1},\\dài \tamthoi{170 m}, ở hai nhánh\\Lê Văn Lương};
  \node[red, anchor=west, align=left] at (1.2,-1.8) {xe rẽ trái trái phép từ Hoàng Đạo Thúy:\\vào nút, dừng chờ, rồi rẽ};
  \node[anchor=west, align=left] at (1.4,2.2) {vùng chờ giữa nút\\(liên kết nội bộ)};
  \draw[->, gray] (1.35,2.0) -- (0.3,0.35);
\end{tikzpicture}
\caption{Sơ đồ minh họa các bản vá mạng (không theo tỷ lệ). Đường đỏ vẽ xe từ nhánh Hoàng Đạo Thúy phía nam; nhánh phía bắc đối xứng. Chiều dài \tamthoi{30~m} lấy từ khảo sát hiện trường, chưa đo; \tamthoi{170~m} là giả định lấy từ thí nghiệm khả thi. Hình học thật của nút chưa đo.}
\label{fig:3.1}
\end{figure}

\begin{table}[ht]
\centering
\small
\caption{Bốn bản vá mạng cho hành vi cần không gian riêng.}
\label{table:3.0}
\begin{tabular}{|p{3.0cm}|p{11cm}|}
\hline
\textbf{Bản vá} & \textbf{Nội dung} \\
\hline
Rẽ trái trái phép & Trên hai nhánh Hoàng Đạo Thúy (cấm rẽ trái), các kết nối rẽ trái hợp lệ bị xóa rồi thêm lại với quyền dùng chỉ dành cho \texttt{custom1}; chỉ nhóm xe vi phạm thực hiện được chuyển động này. Việc dừng chờ giữa nút dùng điểm chờ nội bộ (internal junction) sẵn có của SUMO, chọn sau thí nghiệm đo số xe máy cùng dừng chờ trên liên kết rẽ trái trong nút (Mục~\ref{section:box_capacity}). Vì xe chờ có thể lệch tới mép làn nội bộ, lượt \texttt{netconvert} cuối đặt \texttt{internal-junctions.vehicle-width} bằng 3,5~m để điểm chờ nằm ngoài dòng đi thẳng đối diện; mọi điểm chờ khác được ghim về vị trí cũ bằng \texttt{contPos} \\
\hline
Làn rẽ phải liên tục & Làn ngoài cùng bên phải (lane 0) của mỗi nhánh, chỉ trong 30~m cuối trước vạch dừng; phía trước đoạn này cả ba làn là làn thường. Liên kết rẽ phải mang trạng thái đèn \texttt{g} (xanh, phải nhường) ở mọi pha. Kết nối đi thẳng từ lane 0 mở cho mọi vClass và theo đèn đi thẳng của nhánh; ở hai nhánh Lê Văn Lương liên kết này mang \texttt{g} vì phải nhường dòng từ vỉa hè đổ vào cùng làn ra. Rẽ trái và quay đầu chỉ có kết nối từ lane 2. Mạng 29/10 có thêm các kết nối chỉ cho \texttt{custom1}, \texttt{custom2}: rẽ trái từ lane 0 và quay đầu từ lane 0, lane 1. Các kết nối này bị bỏ ngày 04/11/2026; xe vi phạm đi các làn đó được mô hình hóa lại bằng lựa chọn cắt về của luật \texttt{queue\_space\_seeking} (Mục~\ref{section:3.3}) \\
\hline
Đoạn làn rẽ phải liên tục & Mỗi nhánh tách thành đoạn phía trước và một đoạn dài \tamthoi{30~m} sát vạch dừng (\texttt{right\_lane\_length\_m}). Đầu đoạn này cũng là điểm M8 đọc làn của xe (Mục~\ref{section:3.7}). Không có cọc phân cách: xe đổi làn tự do giữa lane 0 và lane 1 trên cả hai đoạn. Xe có ý định đi thẳng, rẽ trái hoặc quay đầu mà dừng chờ đèn đỏ ở lane 0 trong đoạn này bị tính là chặn làn rẽ phải \\
\hline
Làn vỉa hè & Mỗi nhánh Lê Văn Lương có một làn vỉa hè chỉ cho \texttt{custom1}, bắt đầu \tamthoi{170~m} trước vạch dừng và chạy tới vạch dừng. Mỗi làn có hai liên kết đèn riêng (24, 25 ở nhánh tây; 26, 27 ở nhánh đông), được ưu tiên, và xung đột với dòng rẽ phải từ lane 0. Xe rẽ phải từ vỉa hè nhường dòng đi thẳng Hoàng Đạo Thúy; một \texttt{prohibition} bỏ chiều nhường ngược lại, vì hai dòng nhường nhau cả hai chiều đã gây kẹt cứng ở cửa ra vỉa hè \\
\hline
\end{tabular}
\end{table}

Chiều dài \tamthoi{30~m} của đoạn làn rẽ phải lấy từ khảo sát hiện trường, chưa đo. Vị trí, bề rộng và tốc độ trên vỉa hè là giá trị giả định lấy từ thí nghiệm khả thi Phase 3B (Mục~\ref{subsection:phase3}); các trường này trong tệp cấu hình hiện trường đo thực vẫn để trống.

\subsection{Thông số hiện trường}

Chưa có số liệu đo thực địa. Tệp cấu hình hiện trường chỉ ghi các điều quan sát được tại nút (nhánh cấm rẽ trái, làn rẽ phải liên tục, đoạn 30~m, các làn mà xe vi phạm dùng để rẽ trái và quay đầu); mọi giá trị đo khác để trống. Mô phỏng dùng một cấu hình đại diện, trong đó mỗi giá trị gắn nguồn và mức bằng chứng (lấy từ tài liệu, chuyển giao từ bối cảnh khác, hoặc giả định).

\textbf{Quy ước giá trị tạm thời.} Số in \tamthoi{màu đỏ} trong luận văn là giá trị tạm thời. Đó là giá trị của cấu hình hiện trường đại diện, giá trị mục tiêu của dữ liệu đại diện (Bảng~\ref{table:3.targets}), tham số hành vi giả định, và các ngưỡng suy từ những giá trị này. Chúng được chọn để chạy thử và kiểm chứng toàn bộ quy trình khi chưa có số đo tại nút. Khi có số đo hiện trường và dữ liệu video, từng giá trị sẽ được ghi lại, thay bằng giá trị thật, và các thực nghiệm được chạy lại. Kết quả mô phỏng ở Chương~\ref{chapter:Experiment} đều tính trên các giá trị tạm thời này, nên không tô màu riêng.

Về làn đường, mọi làn vào và làn ra rộng \tamthoi{3,5~m}, theo mục 9.2.3, Bảng 10 của TCVN 13592:2022 cho đường phố chính~\cite{tcvn13592}; số làn lấy từ OpenStreetMap. Mỗi nhánh có ba làn. Thượng nguồn cả ba là làn thường; ở \tamthoi{30~m} cuối, lane 0 là làn rẽ phải liên tục (xác nhận khi khảo sát hiện trường), trong đó rẽ phải không theo đèn và đi thẳng theo đèn đi thẳng. Lane 1 đi thẳng; lane 2 đi thẳng và quay đầu, thêm rẽ trái ở hai nhánh Lê Văn Lương. Cách chia hướng theo làn là giả định.

Về đèn tín hiệu, chu kỳ dài \tamthoi{120~s}. TCCS 24:2018/TCĐBVN~\cite{tccs24_2018} (mục 6.7.5) cho chu kỳ thường 60--90~s và tới 120~s trong trường hợp đặc biệt; Mai và cộng sự đo được 127~s tại nút Láng Hạ -- Láng trên cùng trục~\cite{mai2019signal}, nên luận văn lấy cận trên 120~s. Theo khảo sát hiện trường, đèn có ba pha: Lê Văn Lương đi thẳng; mũi tên rẽ trái và quay đầu của Lê Văn Lương (hai nhánh cùng lúc, đi thẳng đỏ); và Hoàng Đạo Thúy. Vàng \tamthoi{4~s} (TCCS 24:2018, mục 6.7.6, với tốc độ thiết kế không quá 60~km/h), đỏ toàn phần \tamthoi{2~s} (giả định). Chia xanh \tamthoi{42/16/44~s} là giả định: 64~s xanh của Lê Văn Lương ở phiên bản hai pha được tách thành 42~s đi thẳng, 16~s mũi tên và 6~s chuyển pha. Các trạng thái của chu kỳ là \tamthoi{42/4/2/16/4/2/44/4/2~s}.

Về hướng đi và không gian của xe vi phạm, hai nhánh Hoàng Đạo Thúy cấm rẽ trái, còn rẽ trái từ Lê Văn Lương được phép và có mũi tên riêng (khảo sát hiện trường). Việc xe vi phạm rẽ trái và chờ giữa nút ở hai nhánh cấm là giả định. Xe vi phạm rẽ trái từ lane 0 và quay đầu từ lane 0, lane 1 (khóa \texttt{violator\_turn\_lanes}) là quan sát của tác giả ngày 29/10/2026, chưa đếm. Hai nhánh Lê Văn Lương có vỉa hè mà xe máy dùng để vượt hàng chờ (khảo sát hiện trường); hình học vỉa hè là giả định.

Khi có số liệu thực địa, các giá trị này sẽ được thay. Luận văn vì vậy không khẳng định mạng mô phỏng có đúng hình học hay chương trình đèn của nút thật.

\section{Biểu diễn luật hành vi}
\label{section:3.3}

\subsection{Đặc tả luật}

Một đặc tả luật có phần định danh, phần truy vết và hai thành phần chính (Bảng~\ref{table:3.1a}).

\begin{table}[ht]
\centering
\small
\caption{Cấu trúc của một đặc tả luật.}
\label{table:3.1a}
\begin{tabular}{|p{3.4cm}|p{10.6cm}|}
\hline
\textbf{Thành phần} & \textbf{Nội dung} \\
\hline
Định danh, nguồn gốc & Mã đặc tả; nguồn gốc là thủ công, LLM hoặc Optuna \\
\hline
Truy vết & Mô hình, phiên bản prompt, vòng lặp, chỉ số mẫu \\
\hline
Nguyên mẫu tài xế (driver archetype) & Danh sách các nguyên mẫu. Mỗi nguyên mẫu có: vClass cơ sở; phân phối kiểu phương tiện (\texttt{vTypeDistribution}) mà nó thuộc về; tỷ trọng trong phân phối đó; tập tham số vType; lý do chọn (bắt buộc với đặc tả do LLM sinh). Tổng tỷ trọng trong mỗi phân phối bằng 1 \\
\hline
Luật TraCI có điều kiện & Danh sách luật. Mỗi luật có loại luật, các nguyên mẫu áp dụng và tham số \\
\hline
\end{tabular}
\end{table}

Mặc định mỗi vClass cơ sở thuộc một phân phối riêng. Nguyên mẫu \texttt{custom1} và \texttt{custom2} có thể khai báo khóa \texttt{distribution} để nằm trong phân phối xe máy hoặc ô tô tuân thủ. Xe của nguyên mẫu đó đi các hướng hợp lệ của nhu cầu nhưng được dùng không gian dành cho xe vi phạm, ví dụ vỉa hè.

\subsection{Tham số tĩnh}

Không gian tham số tĩnh nằm trong một tệp miền giá trị dùng chung cho LLM và Optuna. Mỗi dòng ghi: mô hình con của SUMO, đơn vị, cận dưới, cận trên, giá trị mặc định, nhóm hành vi liên quan, các chỉ số có thể ràng buộc tham số, và đường dẫn tài liệu SUMO xác nhận tên thuộc tính. Bảng~\ref{table:3.1b} liệt kê các nhóm.

\begin{table}[ht]
\centering
\small
\caption{Các nhóm tham số tĩnh trong không gian tìm kiếm.}
\label{table:3.1b}
\begin{tabular}{|p{4.2cm}|p{9.8cm}|}
\hline
\textbf{Nhóm} & \textbf{Thuộc tính vType} \\
\hline
Bám xe (Krauss) & \texttt{tau}, \texttt{minGap}, \texttt{accel}, \texttt{decel}, \texttt{sigma}, \texttt{speedFactor}, \texttt{speedDev} \\
\hline
Làn con & \texttt{minGapLat}, \texttt{maxSpeedLat}, \texttt{latAlignment} \\
\hline
Chuyển làn SL2015 & \texttt{lcSublane}, \texttt{lcPushy}, \texttt{lcPushyGap}, \texttt{lcAssertive}, \texttt{lcImpatience}, \texttt{lcTimeToImpatience}, \texttt{lcMaxSpeedLatStanding}, \texttt{lcSpeedGain} \\
\hline
Nút giao & \texttt{impatience}, \texttt{jmDriveAfterRedTime}, \texttt{jmDriveRedSpeed}; thêm \texttt{jmTimegapMinor} chỉ cho \texttt{custom1} \\
\hline
\end{tabular}
\end{table}

Hai thuộc tính không chỉ số nào ràng buộc (\texttt{lcAccelLat}, \texttt{jmDriveAfterYellowTime}) được cố định ở giá trị mặc định. \texttt{jmIgnoreFoeProb} và \texttt{jmIgnoreFoeSpeed} bị loại khỏi không gian tìm kiếm.

\subsection{Luật có điều kiện}

Bảng~\ref{table:3.2} định nghĩa sáu loại luật TraCI. Miền giá trị của tham số là lựa chọn mô hình hóa, phần lớn ở mức giả định, và do bộ kiểm tra hợp lệ kiểm soát.

Luật \texttt{queue\_space\_seeking} gộp hành vi 5, 6, 7 của nhóm F2 vào một cơ chế và thay luật đi sai làn riêng lẻ của phiên bản trước. Luật không có lựa chọn đi ngược chiều vì hành vi 8 đã hoãn. Xác suất của một lựa chọn không có trên mạng được chia lại theo tỷ lệ cho các lựa chọn còn lại. Trên mạng 04/11, lựa chọn \texttt{right\_lane} không tồn tại, vì nó chỉ có nghĩa khi lane 0 cấm đi thẳng và có cọc ngăn đổi làn; hành vi 6 đến từ cách SUMO chọn làn và từ hàng chờ. Lựa chọn \texttt{sidewalk} áp dụng ở hai nhánh Lê Văn Lương.

Lựa chọn \texttt{wrong\_lane} là kiểu cắt về (cut-back), dùng từ ngày 08/11/2026. Xe rẽ trái chuyển sang lane 0; xe quay đầu chuyển sang lane 0 hoặc lane 1, chọn làn có hàng chờ ngắn hơn. Các làn này lấy từ khóa \texttt{violator\_turn\_lanes} của cấu hình hiện trường và không có kết nối tới hướng của xe. Khi xe vào đoạn \tamthoi{30~m} sát vạch dừng, luật ra lệnh đổi về lane 2, với bit 9--8 của lane change mode bằng 01 và \texttt{lcPushy} bằng 1 như luật \texttt{uturn\_bypass}; hai giá trị này được trả lại khi xe qua vạch. Chưa nhập được lane 2 thì xe đứng ở cuối làn sai và chặn xe sau. Lựa chọn bị bỏ qua cho hướng bị cấm (rẽ trái từ hai nhánh Hoàng Đạo Thúy) và cho xe đã ở trong đoạn 30~m. Trước ngày 08/11, lựa chọn này chỉ có một dạng: xe vi phạm quay đầu chuyển từ lane 2 sang lane 1. Ngoài lúc cắt về, luật không khóa chế độ chuyển làn của xe. Mỗi can thiệp ghi lựa chọn vào nhật ký. Bộ kiểm tra hợp lệ và bộ điều khiển dùng chung một mô-đun để xác định lựa chọn đi sai làn nào dùng được trên mạng.

Luật \texttt{swarm\_red\_running} có thêm điều kiện nút thông từ ngày 04/11/2026, sau chẩn đoán khóa nút của B1 trên mạng 29/10 (Mục~\ref{subsection:gridlock}). Xe đã rút ``vượt đỏ'' chỉ vào nút khi làn ra của liên kết còn chỗ (chiều dài xe cộng \texttt{minGap}), không có xe trên các làn nội bộ cắt liên kết, và không có xe đứng trên làn nội bộ của chính liên kết. Nút chưa thông thì xe đợi ở vạch dừng. Nếu đèn hết đỏ trước khi nút thông, lượt đó được ghi là không thực hiện (\texttt{junction\_blocked}). Điều kiện này không thêm tham số.

Luật \texttt{uturn\_bypass} được thêm sau khảo sát hiện trường ngày 30/10/2026: phần lớn xe máy đứng sau một xe quay đầu đang chờ sẽ lách qua, ô tô ít khi làm vậy. Từ ngày 04/11/2026 luật thuộc cấu trúc cố định của B1, với $p = \tamthoi{0{,}8}$ cho xe máy, $p = \tamthoi{0{,}1}$ cho ô tô và $t_{\text{wait}} = \tamthoi{2}$~s, giả định từ mô tả định tính (Mục~\ref{subsection:uturn}), và có trong prompt từ \texttt{v003}. Xe quay đầu đứng ở lane 0 hoặc lane 1, nơi không có kết nối quay đầu, cũng được tính là xe chặn. Trong lúc can thiệp, bit 9--8 của lane change mode đặt thành 01 (chỉ tránh va chạm tức thời) và \texttt{lcPushy} của riêng xe đó đặt thành 1; cả hai được trả lại khi can thiệp kết thúc. Speed mode không đổi.

\begin{table}[ht]
\centering
\small
\caption{Các luật TraCI có điều kiện.}
\label{table:3.2}
\begin{tabular}{|p{2.6cm}|p{5.2cm}|p{3.3cm}|p{2.1cm}|}
\hline
\textbf{Luật} & \textbf{Hành vi} & \textbf{Tham số (miền)} & \textbf{vClass} \\ \hline
\texttt{early\_\allowbreak start} & Trong $t_{\text{before\_green}}$ giây cuối pha đỏ, xe đầu hàng nhích qua vạch dừng với xác suất $p$, tốc độ không quá $v_{\max}$ & $t_{\text{before\_green}} \in [0,10]$~s; $p \in [0,1]$; $v_{\max} \in [0{,}5;\,5]$~m/s & passenger, motorcycle, custom1 \\ \hline
\texttt{swarm\_\allowbreak red\_\allowbreak running} & Xác suất vượt đỏ $p = \min(p_0 + k\,n,\ p_{\max})$, với $n$ là số xe đã vượt đỏ trên cùng nhánh trong $\Delta t$ giây gần nhất. Xe đã rút vượt đỏ chỉ vào nút khi nút thông & $p_0, p_{\max} \in [0,1]$; $k \in [0;\,0{,}5]$~1/xe; $\Delta t \in [1,30]$~s & passenger, motorcycle, custom1 \\ \hline
\texttt{queue\_\allowbreak space\_\allowbreak seeking} & Mỗi xe rút ngẫu nhiên một lần khi điều kiện kích hoạt đầu tiên thỏa: hàng chờ ở làn đúng dài ít nhất $q$; hoặc xe đã dừng ít nhất $t_w$; hoặc, $t_{g}$ giây trước lượt xanh kế tiếp, xe ước lượng (với sai số nhân ngẫu nhiên $\sigma$) rằng mình không qua kịp vạch trong lượt xanh đó. Kết quả rút là một trong \texttt{wrong\_lane} (đi làn sai rồi cắt về lane 2 ở đoạn 30~m cuối), \texttt{right\_lane}, \texttt{sidewalk} hoặc không làm gì; xe đi trong không gian thay thế với tốc độ không quá $v_{\text{creep}}$ & $q \in [5,80]$~m; $t_w \in [5,120]$~s (tùy chọn); $t_{g} \in [0,20]$~s; $\sigma \in [0,1]$; ba xác suất trong $[0,1]$, tổng $\le 1$; $v_{\text{creep}} \in [0{,}5;\,5]$~m/s; vùng nhập lại $[3,20]$~m & custom1, custom2 \\ \hline
\texttt{illegal\_\allowbreak left\_\allowbreak box} & Giữ thêm xe đang chờ giữa nút một khoảng $\Delta_{\text{release}}$ so với lúc SUMO tự giải phóng; không bỏ qua xe xung đột & $\Delta_{\text{release}} \in [0,4]$~s & custom1 \\ \hline
\texttt{slow\_\allowbreak lane\_\allowbreak center} & Xe máy đi chậm, giữ vị trí ngang lệch $\delta$ so với tim làn & $\delta \in [-1,1]$~m & motorcycle \\ \hline
\texttt{uturn\_\allowbreak bypass} & Xe đứng sau một xe quay đầu đang đứng trên cùng làn vạch dừng (cách không quá 10~m, chồng ngang với nó) khi đèn của chính nó xanh. Sau $t_{\text{wait}}$ giây, rút một lần với xác suất $p$ cho mỗi cặp xe. Nếu trúng, xe lách ngang vào khe gần nhất bên cạnh xe quay đầu (lề 0,1~m); không có khe thì đổi sang làn kề có nối tới hướng đi của xe & $p \in [0,1]$; $t_{\text{wait}} \in [0,10]$~s & passenger, motorcycle, custom1, custom2 \\ \hline
\end{tabular}
\end{table}

\section{Kiểm tra hợp lệ và biên dịch}
\label{section:3.4}

Bộ kiểm tra hợp lệ chạy năm nhóm kiểm tra. Nhóm (i) đối chiếu đặc tả với lược đồ JSON. Nhóm (ii) kiểm tra tên thuộc tính có trong miền tham số, giá trị nằm trong khoảng cho phép, tổng tỷ trọng theo từng phân phối, và tham số của từng loại luật. Nhóm (iii) tìm mâu thuẫn logic: \texttt{minGapLat} lớn hơn bề rộng làn nhỏ nhất; khai báo cả \texttt{lcPushy} và \texttt{lcPushyGap} (SUMO bỏ qua \texttt{lcPushyGap}); \texttt{lcPushyGap} lớn hơn \texttt{minGapLat}; $p_0 > p_{\max}$; luật áp dụng cho vClass không phù hợp; lựa chọn đi sai làn trên nhánh không có kết nối tương ứng; mã nguyên mẫu trùng tên kiểu xe mặc định của SUMO (\texttt{DEFAULT\_*}). Nhóm (iv) chỉ chạy ở chế độ P1 và yêu cầu cấu trúc khớp B1 (Mục~\ref{section:3.8}). Nhóm (v) chỉ chạy ở chế độ P2, từ prompt \texttt{v005}: đặc tả phải dùng đủ sáu loại luật TraCI và mỗi loại phải được bật ở ít nhất một mục, tức $p > 0$ với \texttt{early\_start} và \texttt{uturn\_bypass}, $p_0 > 0$ với \texttt{swarm\_red\_running}, và ít nhất một xác suất lựa chọn lớn hơn 0 với \texttt{queue\_space\_seeking}; \texttt{illegal\_left\_box} và \texttt{slow\_lane\_center} không có xác suất nên chỉ cần có mặt.

Kết quả là danh sách lỗi có cấu trúc, phân theo mã: \texttt{schema}, \texttt{unknown\_attr}, \texttt{out\_of\_bounds}, \texttt{share\_sum}, \texttt{rule\_params}, \texttt{logic}, \texttt{structure\_locked} (nhóm iv) và \texttt{coverage} (nhóm v). Nhóm (v) được thêm ngày 08/11/2026, sau khi các đặc tả P2 sinh với prompt \texttt{v004} bỏ \texttt{queue\_space\_seeking}, \texttt{uturn\_bypass} và \texttt{slow\_lane\_center}, tức bỏ một phần các hành vi của kịch bản (Mục~\ref{subsection:campaigns_n1108}). Chiến dịch dùng \texttt{v004} trở về trước không chịu kiểm tra này, nên nhật ký cũ vẫn phát lại được. Danh sách này gửi ngược lại cho LLM được. Tỷ lệ đặc tả hợp lệ được ghi cho từng lần LLM trả lời và cho trạng thái cuối của từng mẫu, và báo cáo như một chỉ số của vòng lặp LLM.

Bộ biên dịch chuyển đặc tả hợp lệ thành hai tệp. \texttt{vtypes.add.xml} chứa các phân phối kiểu phương tiện \texttt{car}, \texttt{moto}, \texttt{moto\_violator} và, nếu đặc tả có nguyên mẫu riêng, \texttt{car\_violator}; cả bốn dùng mô hình chuyển làn SL2015. \texttt{traci\_rules.json} chứa các luật có điều kiện. Tên phân phối khớp với tên trong tệp nhu cầu, nên mọi đặc tả dùng chung một nhu cầu.

\section{Bộ điều khiển thực thi luật}
\label{section:3.5}

Mỗi loại luật là một lớp riêng, được gọi ở mỗi bước mô phỏng. Bộ điều khiển dùng đăng ký nhận dữ liệu (subscription) thay vì truy vấn từng xe, và chỉ xử lý xe trong bán kính 220~m quanh nút, đủ phủ điểm bắt đầu làn vỉa hè. Trạng thái đèn của mỗi xe lấy theo liên kết TLS của chính nó, nên xe rẽ trái Lê Văn Lương trong pha mũi tên không bị tính là vượt đỏ. Mỗi xe chịu tối đa một can thiệp tại một thời điểm, để các lệnh ghi đè không chồng lên nhau. Luật vượt đèn đỏ chỉ xóa bit ``phanh để không vượt đèn đỏ'' của speed mode; các bit an toàn khác, kể cả bit nhường đường cho xe xung đột, được giữ nguyên.

Sau can thiệp, bộ điều khiển khôi phục speed mode ban đầu và kiểm tra lại giá trị SUMO trả về. Mỗi can thiệp được ghi vào nhật ký: thời điểm, mã xe, luật, hành động, việc SUMO có thực thi hay không, lý do nếu không, và thời điểm kết thúc. Số va chạm liên quan đến xe đang chịu can thiệp được thống kê riêng trong \texttt{manifest.json} của mỗi lần chạy.

\section{Dữ liệu mục tiêu}
\label{section:3.6}

Luận văn dự kiến dùng dữ liệu quan sát trích từ video quay tại nút. Lược đồ dữ liệu quan sát đã được định nghĩa (Phụ lục~\ref{section:appendix_schema}). Nó gồm các bảng số đếm rẽ theo 15 phút, sự kiện đèn, lượt qua vạch dừng, ảnh chụp hàng chờ, vị trí ngang, hàng chờ giữa nút và lượt đi qua nút, cộng ba bảng cho nhóm F2 là chuỗi thời gian hàng chờ, sự kiện tìm khoảng trống và độ trễ của xe rẽ phải. Đầu ra mô phỏng được chuyển sang cùng lược đồ, nên một bộ hàm chỉ số tính được cho cả hai nguồn.

\begin{sloppypar}
Chưa có dữ liệu video \textbf{[TODO: cập nhật khi có dữ liệu thực]}. Để kiểm thử cả quy trình, luận văn dựng năm ``ngày đại diện'' PROXY\_D1--PROXY\_D5, sinh theo từng sự kiện phương tiện từ một bảng giá trị mục tiêu tổng hợp từ tài liệu. Mỗi giá trị mục tiêu mang nhãn bằng chứng: lấy từ tài liệu, chuyển giao từ bối cảnh khác, hoặc giả định. Biến động giữa các ngày đến từ hệ số theo ngày, chuẩn hóa để trung bình năm ngày bằng giá trị mục tiêu danh nghĩa.
\end{sloppypar}

Bảng~\ref{table:3.targets} liệt kê giá trị mục tiêu danh nghĩa chính của từng nhóm chỉ số. Mọi giá trị trong bảng là giá trị tạm thời và sẽ được thay bằng số đo từ video. Tệp mục tiêu còn các tham số hình dạng (họ phân phối, hệ số biến thiên, hệ số theo ngày); phần lớn chúng mang nhãn giả định và cũng là giá trị tạm thời.

\begin{table}[ht]
\centering
\small\setlength{\tabcolsep}{3pt}
\caption{Giá trị mục tiêu danh nghĩa chính của dữ liệu đại diện (giá trị tạm thời). LVL: hai nhánh Lê Văn Lương; HDT: hai nhánh Hoàng Đạo Thúy. Nhãn: CG là chuyển giao từ bối cảnh khác, GĐ là giả định.}
\label{table:3.targets}
\begin{tabular}{|p{1.0cm}|p{6.2cm}|p{4.4cm}|p{2.4cm}|}
\hline
\textbf{Nhóm} & \textbf{Đại lượng} & \textbf{Giá trị} & \textbf{Nhãn} \\ \hline
G0 & Lưu lượng qua vạch theo nhánh LVL\_W, LVL\_E, HDT\_N, HDT\_S (xe/h) & \tamthoi{2620; 2130; 1800; 1630} & tổng CG~\cite{mai2019signal}, chia nhánh GĐ \\ \hline
G0 & Tỷ lệ xe máy & \tamthoi{0,865} & CG~\cite{mai2019signal} \\ \hline
G0 & Chia hướng đi thẳng/rẽ phải/rẽ trái/quay đầu & \tamthoi{0,70/0,15/0,10/0,05}; nhánh cấm rẽ trái \tamthoi{0,80/0,15/--/0,05} & GĐ \\ \hline
M1 & Khoảng thời gian bão hòa của dòng xe máy (s); hệ số ô tô & \tamthoi{0,254}; \tamthoi{1,4} & CG~\cite{vu2015saturation} \\ \hline
M2 & Tỷ lệ xuất phát sớm, xe máy LVL; HDT & \tamthoi{0,0945}; \tamthoi{0,1582} & CG~\cite{mai2019signal} \\ \hline
M2 & Tỷ lệ xuất phát sớm, ô tô LVL; HDT & \tamthoi{0,0027}; \tamthoi{0,0068} & CG~\cite{mai2019signal} \\ \hline
M2 & Độ trễ từ lúc xanh đến xe đầu tiên qua vạch (s) & \tamthoi{1,2} & GĐ \\ \hline
M3 & Tỷ lệ vượt đỏ, xe máy LVL; HDT & \tamthoi{0,0649}; \tamthoi{0,1106} & CG~\cite{mai2019signal} \\ \hline
M3 & Tỷ lệ vượt đỏ, ô tô LVL; HDT & \tamthoi{0,0027}; \tamthoi{0,0050} & CG~\cite{mai2019signal} \\ \hline
M4 & Tỷ lệ xe máy dừng chờ vượt quá vạch dừng; độ sâu (m) & \tamthoi{0,63}; \tamthoi{0--4} & CG~\cite{zulkiffli2023stopping}; GĐ \\ \hline
M5 & Diện tích hàng chờ của một xe máy; một ô tô (m\textsuperscript{2}) & \tamthoi{2,6}; \tamthoi{17,4} & GĐ \\ \hline
M6 & Tốc độ tiếp cận; tốc độ qua nút của xe máy (km/h) & \tamthoi{27,3}; \tamthoi{21,3} & CG~\cite{chu2005speed} \\ \hline
M6 & Tốc độ tiếp cận; tốc độ qua nút của ô tô (km/h) & \tamthoi{25}; \tamthoi{20} & GĐ \\ \hline
M7 & Số xe chờ giữa nút mỗi chu kỳ, mỗi nhánh có xe rẽ trái trái phép & \tamthoi{6} & GĐ \\ \hline
M7 & Độ trễ giải phóng vùng chờ (s); tổn thất khởi động của hướng giao cắt (s/xe) & \tamthoi{5,65}; \tamthoi{0,3} & GĐ (dựa trên~\cite{hsu2024v2x}); GĐ \\ \hline
M8 & Tỷ lệ xe rẽ đi sai làn, xe máy; ô tô & \tamthoi{0,25}; \tamthoi{0,05} & GĐ \\ \hline
M9 & Xác suất xe đi thẳng chọn lane 0, xe máy; ô tô & \tamthoi{0,05}; \tamthoi{0,02} & GĐ \\ \hline
M9 & Thời gian xe đi thẳng chiếm lane 0 (s) & \tamthoi{12} & GĐ \\ \hline
M10 & Xác suất xe máy đi thẳng chọn vỉa hè & \tamthoi{0,03} & GĐ \\ \hline
\end{tabular}
\end{table}

Từ ngày 08/11/2026, M8 đọc làn của xe tại đầu làn rẽ phải liên tục, cách vạch dừng \tamthoi{30~m} (Mục~\ref{section:3.7}). Khi đo bằng video, làn của xe phải được đọc tại đúng điểm này. Bộ sinh dữ liệu đại diện chỉ có một vị trí ngang cho mỗi xe, nên mục tiêu M8 không đổi khi định nghĩa đổi.

Việc chia ngày như sau. PROXY\_D1--PROXY\_D3 dùng để hiệu chỉnh. PROXY\_D4--PROXY\_D5 chỉ dùng cho đánh giá cuối; LLM và Optuna không được thấy thống kê của chúng. Các mục tiêu của nhóm F2 (xác suất chọn không gian thay thế, độ trễ xe rẽ phải bị chặn) đều mang nhãn giả định. Hai mục tiêu của hành vi 6 (\texttt{right\_lane\_choice\_prob}, \texttt{lane0\_hold\_s}) được suy từ thiết lập có cọc của thí nghiệm 3A và chưa định nghĩa lại theo mô hình làn 30~m, nên điểm M9 trên mạng 29/10 chưa đọc được như sai lệch hành vi (Mục~\ref{subsection:ref_n29}). Bộ năm ngày hiện tại được sinh lại cho kịch bản ba pha; ngày dữ liệu không được chọn lại theo kết quả. Dữ liệu mục tiêu không phải quan sát thực, nên kết quả hiệu chỉnh trên các ngày này chỉ cho biết quy trình vận hành được, không cho biết mô phỏng giống giao thông thực đến đâu.

\section{Hệ thống chỉ số đánh giá}
\label{section:3.7}

\subsection{Các chỉ số}

Mỗi chỉ số là một hàm trên dữ liệu theo lược đồ chung, trả về một phân phối hoặc một giá trị vô hướng. Bảng~\ref{table:3.3} liệt kê các chỉ số đã triển khai. M1--M10 được tính trong cửa sổ $[300, 1800)$~s ở cả phía mô phỏng và phía quan sát (300~s đầu là giai đoạn khởi động). Cổng G0 tính trên hai ô 15 phút của khoảng 0--1800~s. M1--M3 chỉ xét các hướng đi theo đầu đèn đi thẳng của nhánh: rẽ phải trên làn liên tục và các hướng có mũi tên riêng (rẽ trái, quay đầu của Lê Văn Lương) không được phân loại, ở cả phía mô phỏng và phía quan sát.

\begin{table}[ht]
\centering
\small
\caption{Các chỉ số đánh giá.}
\label{table:3.3}
\begin{tabular}{|p{1.0cm}|p{9.0cm}|p{3.2cm}|}
\hline
\textbf{Mã} & \textbf{Định nghĩa} & \textbf{Đơn vị / dạng} \\ \hline
G0 & Cổng lọc: GEH $=\sqrt{2(m-c)^2/(m+c)}$ của số đếm rẽ theo ô 15 phút $\times$ nhánh $\times$ hướng $\times$ vClass ($m$: mô phỏng, $c$: quan sát). Lần chạy đạt khi tỷ lệ ô có GEH $<5$ không nhỏ hơn 0,85. Không tính vào điểm & Tỷ lệ ô đạt \\ \hline
M1 & Khoảng thời gian xả hàng (discharge headway) sau khi đèn xanh, theo vClass & s; phân phối \\ \hline
M2 & Tỷ lệ xe qua vạch trong 10~s cuối pha đỏ; độ trễ từ lúc xanh đến khi xe đầu tiên qua vạch & Tỷ lệ; s \\ \hline
M3 & Số xe vượt đỏ mỗi chu kỳ, tỷ số phương sai trên trung bình, tỷ lệ vượt đỏ & xe/chu kỳ; tỷ lệ \\ \hline
M4 & Vị trí ngang của xe máy đang dừng khi đèn chuyển xanh; tỷ lệ và độ sâu vượt quá vạch dừng & m; phân phối \\ \hline
M5 & Chiều dài hàng chờ và mật độ xe máy khi đèn chuyển xanh, theo nhánh & m; xe/m\textsuperscript{2} \\ \hline
M6 & Thời gian đi qua nút và qua nhánh vào, theo hướng $\times$ vClass & s; phân phối \\ \hline
M7 & M7a: số xe chờ giữa nút và thời điểm xả. M7b: độ trễ giải phóng vùng chờ (đại diện cho PET). M7c: độ dốc tổn thất khởi động của hướng giao cắt theo tổng số xe chờ ở hai vùng chờ & xe; s; s/xe \\ \hline
M8 & Đi sai làn: tỷ lệ xe rẽ (phải, trái, quay đầu) đứng ở làn không khớp hướng ra. Làn đọc ở lần quan sát cuối trước khi đầu xe tới đầu làn rẽ phải liên tục, \tamthoi{30~m} trước vạch dừng (trước 08/11/2026: tại vạch dừng) & Tỷ lệ \\ \hline
M9 & Chặn làn rẽ phải liên tục. M9a: số xe có hướng đi thẳng, rẽ trái hoặc quay đầu đứng chờ đèn đỏ của chính nó trên lane 0 trong đoạn 30~m cuối, mỗi chu kỳ. M9b: độ trễ và tỷ lệ xe rẽ phải bị chặn. M9c: độ dài hàng chờ lúc xe đó đứng vào lane 0 (chỉ báo cáo) & xe/chu kỳ; s; tỷ lệ; m \\ \hline
M10 & Đi trên vỉa hè. M10a: số xe máy trên vỉa hè mỗi chu kỳ. M10b: độ dài hàng chờ lúc xe chọn vỉa hè (chỉ báo cáo) & xe/chu kỳ; m \\ \hline
\end{tabular}
\end{table}

M8 đọc làn ở đầu đoạn 30~m vì xe đi sai làn kiểu cắt về đã nhập lại lane 2 trước khi tới vạch dừng. Nếu đọc tại vạch, M8 không ghi nhận xe đó là đi sai làn (Mục~\ref{subsection:cutback}). Trên nhánh không có đoạn làn rẽ phải, M8 đọc làn tại vạch dừng. Xe xuất hiện lần đầu khi đã ở trong đoạn 30~m bị bỏ khỏi M8.

\subsection{Hàm khoảng cách và điểm tổng}

Với chỉ số dạng phân phối, khoảng cách là khoảng cách Wasserstein bậc 1 chia cho khoảng tứ phân vị (IQR) của dữ liệu quan sát; nếu IQR bằng 0, mẫu số lùi lần lượt về độ lệch chuẩn, trị tuyệt đối của trung bình, rồi 1. Với chỉ số vô hướng, khoảng cách là sai số tương đối. Mọi khoảng cách theo ngày bị chặn trên ở $D_{\text{empty}} = 2$, và khi một phía không có sự kiện hoặc giá trị không xác định, khoảng cách bằng đúng $D_{\text{empty}}$. Quy tắc này ngăn bộ tối ưu hạ điểm bằng cách triệt tiêu một hành vi đến mức không còn sự kiện.

Để tách sai lệch mô hình khỏi biến động tự nhiên, cùng các khoảng cách đó được tính giữa các ngày quan sát với nhau; kết quả là sàn nhiễu $d_{\text{day-day}}$, với sàn dưới 0,05. Mỗi khóa được báo cáo bằng tỷ số $R = d_{\text{sim-obs}}/d_{\text{day-day}}$.

Điểm tổng $S$ dùng cho bộ tối ưu là trung bình có trọng số của các khoảng cách. Trọng số chia đều cho các nhóm M1--M10 có mặt; trong mỗi nhóm chia đều cho các chỉ số con (ví dụ M7a, M7b, M7c), rồi cho các khóa. Chỉ khóa được chấm điểm mới có trọng số. $S$ tách thành hai phần cộng lại bằng $S$: $S_{\text{sourced}}$ từ các khóa có mục tiêu có nguồn, và $S_{\text{assumed}}$ từ các khóa có mục tiêu giả định.

M9c và M10b chỉ báo cáo, không chấm điểm. Trên mạng 29/10, M9c là 1,96 với B0 và 2,00 với B1, vì hàng chờ lúc xe đứng vào lane 0 dài hơn mục tiêu (trung bình 47 và 60~m so với 16~m). M10b bằng $D_{\text{empty}}$ ở cả hai vì không xe nào đi vỉa hè.

Một hạn chế đã biết: với tỷ lệ rất nhỏ (ví dụ tỷ lệ xuất phát sớm của ô tô), khoảng cách khi mô phỏng bằng 0 là 1, còn khi mô phỏng gấp hơn hai lần quan sát là 2. Ép hành vi về 0 vì vậy vẫn rẻ hơn. Cách xử lý trường hợp này chưa chốt \textbf{[Câu hỏi cho GVHD: nên làm mượt khoảng cách của tỷ lệ nhỏ (ví dụ cộng hằng số vào mẫu số, hoặc dùng khoảng cách trên số đếm thay vì tỷ lệ), hay giữ nguyên và nêu là hạn chế?]}.

\section{Quy trình hiệu chỉnh và các phương án đối sánh}
\label{section:3.8}

\subsection{Vòng lặp LLM}

\textbf{Giao diện và đầu ra.} LLM được gọi qua giao diện tương thích OpenAI, để mã không gắn với một runtime. Đầu ra có cấu trúc theo lược đồ JSON của đặc tả. Lược đồ gửi cho LLM là một tập con an toàn với bộ giải mã có cấu trúc: thuộc tính được liệt kê rõ từ tệp miền giá trị, không dùng cấu trúc điều kiện. Lược đồ đầy đủ vẫn dùng để kiểm tra hợp lệ. Nếu máy chủ từ chối lược đồ, trình khách lùi về chế độ JSON thông thường.

\textbf{Prompt và phản hồi.} Prompt có phiên bản, khóa theo chiến dịch; các phiên bản cũ giữ nguyên để phát lại chiến dịch cũ (Bảng~\ref{table:3.prompts}). Prompt gồm: vai trò và ràng buộc; bối cảnh nút giao kèm từ điển tham số sinh tự động; và với mỗi vòng trước, một tin nhắn chứa các đặc tả của vòng đó cùng phản hồi. Phản hồi gồm khoảng cách theo từng chỉ số, tỷ lệ ô GEH đạt, số và vị trí teleport kẹt, lỗi từ bộ kiểm tra hợp lệ, và tóm tắt can thiệp TraCI. Phản hồi chỉ dùng mục tiêu của PROXY\_D1--PROXY\_D3; mã báo lỗi nếu thông tin của ngày đánh giá lọt vào. Từ \texttt{v004}, phản hồi của một vòng gộp thành một bảng: mỗi chỉ số có trọng số một dòng, gồm giá trị quan sát và, với mỗi đặc tả, giá trị mô phỏng và $R$. Vòng gần nhất được gửi đầy đủ (JSON của đặc tả và bảng chỉ số); các vòng cũ hơn chỉ còn vài dòng tóm tắt ($S$, tính đủ điều kiện, nguyên mẫu, luật).

\begin{table}[ht]
\centering
\small
\caption{Các phiên bản prompt.}
\label{table:3.prompts}
\begin{tabular}{|p{1.4cm}|p{1.6cm}|p{10.6cm}|}
\hline
\textbf{Bản} & \textbf{Ngày} & \textbf{Thay đổi} \\ \hline
\texttt{v001} & 27/10 & Kịch bản hai pha, mạng F2 \\ \hline
\texttt{v002} & 28/10 & Kịch bản ba pha sau khảo sát hiện trường \\ \hline
\texttt{v003} & 04/11 & Thêm hành vi 11 (\texttt{uturn\_bypass}); tin nhắn thử lại kèm trạng thái các phân phối vType của đặc tả \\ \hline
\texttt{v004} & 08/11 & Rút gọn cho \texttt{qwen3.5:9b}: lược đồ đầu ra chế độ \texttt{full} có một biến thể cho mỗi vClass cơ sở, chỉ chứa tham số và phân phối của lớp đó, \texttt{distribution} bắt buộc; ba câu nhắc về tổng tỷ trọng, tỷ trọng lớn hơn 0 và mã nguyên mẫu không trùng tên phân phối; chỉ liệt kê 42 khóa có tính điểm; khung mẫu ngắn thay cho bản in đầy đủ của lược đồ JSON; phản hồi và lịch sử rút gọn như trên; mô tả hành vi 5 theo kiểu cắt về \\ \hline
\texttt{v005} & 08/11 & Như \texttt{v004}, thêm yêu cầu dùng đủ sáu luật TraCI ở chế độ \texttt{full} (Mục~\ref{section:3.4}) \\ \hline
\end{tabular}
\end{table}

\textbf{Số lần thử và nhịp vòng.} Mỗi mẫu được thử tối đa 3 lần; khi đặc tả không hợp lệ, danh sách lỗi được gửi lại cho LLM, và đặc tả không hợp lệ không tốn lần chạy SUMO. Mỗi vòng sinh 4 mẫu. Sau ba vòng theo kế hoạch, nếu còn ngân sách (do có mẫu không hợp lệ), vòng lặp chạy thêm cho tới khi hết, tối đa 5 vòng; số vòng thực tế được ghi lại.

\textbf{Tái lập.} Mọi yêu cầu và phản hồi được lưu nguyên văn. Chạy lại với cùng tin nhắn và cùng định danh LLM thì phản hồi đã lưu được phát lại, không gọi lại LLM. Kết quả từng mẫu được lưu ngay sau khi đánh giá, nên vòng lặp chạy tiếp được sau khi bị ngắt.

\textbf{Hai chế độ.} Chế độ \texttt{params\_only} (phương án P1) khóa cấu trúc của B1 (các nguyên mẫu, các thuộc tính khai báo, các luật); LLM chỉ đổi giá trị, và mọi thay đổi cấu trúc bị báo lỗi \texttt{structure\_locked}. Chế độ \texttt{full} (phương án P2) cho LLM sinh cả cấu trúc lẫn tham số.

\textbf{Mô hình LLM.} LLM cục bộ là \texttt{qwen3.5:9b} (9 tỷ tham số, lượng tử hóa Q4\_K\_M, tệp mô hình 5,6~GB), chạy bằng Ollama 0.40.11 trên máy chính (GPU NVIDIA GeForce RTX 4060, 8~GB). Mô hình này vừa với bộ nhớ GPU của máy. Tham số lấy mẫu: nhiệt độ 0,7; seed 42, mẫu thứ $j$ dùng seed $42 + j$ để bốn mẫu trong một vòng khác nhau. Trường \texttt{provenance.model} của mỗi đặc tả ghi định danh mô hình đã sinh ra nó. Ollama chạy với cửa sổ ngữ cảnh 32.768 token. Câu trả lời giới hạn 7000 token (\texttt{max\_tokens}); câu trả lời bị cắt (\texttt{finish\_reason} bằng \texttt{length}) tính là một lần thử hỏng. Chế độ suy nghĩ (thinking) của mô hình là một phần cấu hình của chiến dịch và được ghi vào trạng thái vòng lặp; các chiến dịch P1, P2 ngày 08/11 tắt nó (\texttt{reasoning\_effort} bằng \texttt{none}), vì khi bật, phần suy nghĩ chiếm hết giới hạn token trước khi tới JSON (Mục~\ref{subsection:campaigns_n1108}).

\subsection{Các phương án đối sánh}

Bảng~\ref{table:3.4} định nghĩa sáu cấu hình.

\begin{table}[ht]
\centering
\small
\caption{Sáu cấu hình đối sánh.}
\label{table:3.4}
\begin{tabular}{|c|p{12.5cm}|}
\hline
\textbf{Mã} & \textbf{Định nghĩa} \\
\hline
B0 & SUMO với tham số mặc định, không có luật TraCI \\
\hline
B1 & Bộ luật thủ công; tham số tham khảo mô hình hỗn hợp cho Hà Nội của Saval và cộng sự~\cite{saval2023mixed}. B1 có sẵn cấu trúc nhóm F2: nguyên mẫu xe máy và ô tô tìm khoảng trống với tỷ trọng 0,1, và luật \texttt{queue\_space\_seeking} với mọi xác suất bằng 0 gán cho nguyên mẫu xe máy, để các phương án chỉnh tham số trên cấu trúc B1 điều chỉnh được nhóm hành vi này. Từ ngày 04/11/2026, B1 có thêm luật \texttt{uturn\_bypass} cho mọi nguyên mẫu xe máy ($p = \tamthoi{0{,}8}$) và ô tô ($p = \tamthoi{0{,}1}$), $t_{\text{wait}} = \tamthoi{2}$~s. Nguyên mẫu ô tô tìm khoảng trống chỉ gắn với luật này (Phụ lục~\ref{section:appendix_b1}) \\
\hline
B2 & Tham số do Optuna~\cite{akiba2019optuna} tối ưu trên cấu trúc cố định của B1 \\
\hline
P1 & LLM chỉnh tham số trên cùng cấu trúc cố định của B1. Với cấu trúc B1, tập tham số Optuna và P1 được phép chỉnh trùng nhau \\
\hline
P2 & LLM sinh cả cấu trúc lẫn tham số \\
\hline
H & Cấu trúc tốt nhất của P2, tham số do Optuna tối ưu \\
\hline
\end{tabular}
\end{table}

Hai thí nghiệm so sánh: E1 (P1 so với B2, cùng cấu trúc cố định) đánh giá LLM như một bộ tối ưu; E2 (B2, H và P2) đánh giá giá trị của cấu trúc luật do LLM sinh.

Optuna dùng bộ lấy mẫu TPE với seed 0 và 4 thử nghiệm khởi đầu ngẫu nhiên. Không gian tìm kiếm chỉ gồm tham số mà cấu trúc khai báo, và được tham số hóa lại để mọi thử nghiệm đều hợp lệ: \texttt{lcPushyGap} biểu diễn bằng tỷ lệ của \texttt{minGapLat}; $p_{\max}$ bằng $p_0$ cộng một tỷ lệ của $1-p_0$; tỷ trọng các nguyên mẫu lấy trên đơn hình bằng cách bẻ que (stick-breaking).

Mọi phương án dùng chung ngân sách 36 lần chạy SUMO. Mỗi đặc tả được đánh giá bằng 3 lần chạy, mỗi ngày hiệu chỉnh PROXY\_D1, PROXY\_D2, PROXY\_D3 một lần, với seed 1. Ngân sách vì vậy tương ứng 12 đặc tả: 3 vòng $\times$ 4 mẫu với LLM, hoặc 12 thử nghiệm với Optuna. Cách chia 1 seed $\times$ 3 ngày phủ biến động giữa các ngày trong cùng ngân sách. Bộ đếm ngân sách gắn với từng chiến dịch, lưu trên đĩa và trừ dưới khóa tệp, nên chạy lại một chiến dịch không làm vượt ngân sách. Hai chiến dịch B2 trên mạng ngày 28/10 dùng ngân sách 24 lần chạy (8 thử nghiệm), không dùng 36.

\subsection{Cấu hình mô phỏng và tiêu chí đủ điều kiện}

Mọi lần chạy dùng SUMO 1.27.0 với bước thời gian 0,1~s, độ phân giải làn con 0,4~m, mô hình chuyển làn SL2015, tắt định tuyến lại, bật kiểm tra va chạm trong nút. Thời lượng mô phỏng đầy đủ là 3600~s; vòng hiệu chỉnh và đánh giá dùng 1800~s đầu (07:30--08:00). Giai đoạn phát triển dùng seed 1, 2, 3; đánh giá cuối dự kiến dùng ít nhất 20 seed trên các ngày đánh giá. Mỗi lần chạy ghi một \texttt{manifest.json} gồm mã commit, phiên bản SUMO, giá trị băm của đặc tả và mạng, seed và cấu hình.

Một lần chạy \emph{đủ điều kiện} (eligible) khi thỏa cả năm điều kiện: (i) không có teleport nào không do va chạm; (ii) không còn xe chờ chèn vào mạng khi kết thúc; (iii) qua cổng G0; (iv) số va chạm trên làn không quá \tamthoi{4} va chạm/h và số va chạm trong nút không quá \tamthoi{142} va chạm/h; và (v) có tối đa 1 va chạm liên quan đến xe đang chịu can thiệp TraCI. Hai ngưỡng va chạm được suy từ B0 trên mạng F2 và đóng băng trước vòng LLM đầu tiên (Mục~\ref{subsection:r28_f2}); chúng chưa được tính lại trên mạng 29/10, nơi chính B0 vượt ngưỡng va chạm trên làn (Mục~\ref{subsection:ref_n29}). Trên mạng 04/11, B0 đạt cả năm điều kiện ở 15/15 lần chạy với hai ngưỡng này (Mục~\ref{subsection:ref_n1104}). Hai ngưỡng suy từ B0 trên dữ liệu đại diện nên cũng là giá trị tạm thời. Một đặc tả đủ điều kiện khi mọi lần chạy của nó đủ điều kiện. Nó chỉ được chọn làm tốt nhất nếu tỷ lệ va chạm trên làn trung bình không vượt mức của B0 cộng 1 va chạm/h. Mức của B0 khai báo theo từng chiến dịch: \tamthoi{1,333} va chạm/h trên mạng F2, \tamthoi{2,667} và \tamthoi{1,333} va chạm/h trên hai mạng ngày 28/10, \tamthoi{1,867} va chạm/h trên mạng 04/11 (B0 của \texttt{REF\_n1104s5}). Nếu không có đặc tả nào đủ điều kiện, đặc tả tốt nhất được chọn theo số teleport kẹt và kẹt khi nhường đường trung bình mỗi giờ, rồi theo $S$; khi đó mức trần va chạm không áp dụng.

Việc so sánh các cấu hình trên ngày đánh giá dự kiến dùng kiểm định Mann--Whitney~\cite{mann1947test} cho từng cặp, hiệu chỉnh Holm~\cite{holm1979simple} cho phép so sánh bội, và kích thước hiệu ứng Cliff's delta~\cite{cliff1993dominance}.

Chương này đã trình bày phương pháp theo đúng thứ tự của một vòng hiệu chỉnh. Mạng mô phỏng được sinh hoàn toàn bằng script, với các bản vá tĩnh cho hành vi cần không gian riêng; từ ngày 04/11/2026, xe vi phạm rẽ trái và quay đầu từ làn sai được mô hình hóa bằng luật cắt về thay vì kết nối riêng. Một đặc tả luật gồm các nguyên mẫu tài xế với tham số vType và sáu loại luật TraCI có điều kiện; bộ kiểm tra hợp lệ chặn đặc tả sai trước khi tốn lần chạy, và bộ điều khiển chỉ can thiệp trong giới hạn an toàn. Khi chưa có dữ liệu video, dữ liệu mục tiêu là năm ngày đại diện có nhãn bằng chứng, và mọi giá trị tạm thời được in màu đỏ để thay sau. Mười nhóm chỉ số được gộp thành điểm $S$ có chặn trên, và sáu phương án đối sánh dùng chung ngân sách 36 lần chạy cùng một tiêu chí đủ điều kiện. Chương~\ref{chapter:Experiment} mô tả phần cài đặt và các kết quả thực nghiệm của phương pháp này.

\end{document}
